% GENERATED FILE. Edit canonical lessons, then run npm run generate:books.
% canonical-source-hash: fnv1a64:ee4751857ce02d71
% canonical-lessons: GE-C307-gastgeber, GE-C307-bewirten, GE-C307-ehrengast, GE-C307-ansprache, GE-C307-trinkspruch

\chapter{Der Gastgeber}
\label{ch:ge-der-gastgeber}
\hlchaptermodality{307}

\begin{chapteropening}
\textbf{By the end of this chapter:} \emph{I can talk about hosts and guests: a host, to treat guests, a guest of honour, a speech and a toast.}

Complete the last lesson of chapter 307: I can talk about hosts and guests: a host, to treat guests, a guest of honour, a speech and a toast.
\end{chapteropening}

\section[der Gastgeber]{der Gastgeber — the host}

\label{lesson:GE-C307-gastgeber}



\begin{quote}
Before the new one: say the German for to give presents to, then the German for to unwrap.
\end{quote}

\subsection*{You'll want to know: der Gastgeber}
\textbf{der Gastgeber} — \textquotedblleft{}the host\textquotedblright{}.

\begin{quote}
Der Gastgeber begrüßte jeden Gast an der Tür.
\end{quote}

\begin{quote}
\emph{The host greeted every guest at the door.}
\end{quote}

Plural: \emph{die Gastgeber}. Feminine \emph{die Gastgeberin}.

\begin{cousinweb}
\emph{Gastgeber} joins \emph{Gast} \textquotedblleft{}guest\textquotedblright{} and \emph{Geber} \textquotedblleft{}giver\textquotedblright{}, from \emph{geben} \textquotedblleft{}to give\textquotedblright{}, Proto-Germanic \emph{gebaną} (EtymDB).
\end{cousinweb}

\begin{grammarlens}[title={What you've built}]
One more word for telling, explaining and describing at B1.
\end{grammarlens}

\subsection*{Your turn}
\begin{itemize}
\raggedright
  \item \emph{Say it:} \emph{der Gastgeber}
  \item \emph{Say it:} \emph{der Gastgeber}, once more
  \item \emph{Say it:} a sentence of your own with \emph{der Gastgeber}
  \item \emph{Recall:} say the German for to give presents to, then the German for to unwrap, then say \emph{der Gastgeber} again
\end{itemize}

\begin{tcolorbox}[breakable,colback=teal!4,colframe=teal!35!black,title={Before you move on}]
What does der Gastgeber mean? (\textquotedblleft{}The host\textquotedblright{}.) Say it once more.
\end{tcolorbox}
\section[bewirten]{bewirten — to treat guests, to serve food and drink to}

\label{lesson:GE-C307-bewirten}



\begin{quote}
Before the new one: say the German for to unwrap, then the German for the host.
\end{quote}

\subsection*{You'll want to know: bewirten}
\textbf{bewirten} — \textquotedblleft{}to treat guests, to serve food and drink to\textquotedblright{}.

\begin{quote}
Sie bewirtete ihre Gäste mit Kaffee und Kuchen.
\end{quote}

\begin{quote}
\emph{She treated her guests to coffee and cake.}
\end{quote}

Built on \emph{der Wirt}, the landlord of an inn.

\begin{cousinweb}
\emph{bewirten} is inherited from Middle High German \emph{bewirten} (EtymDB), built on \emph{Wirt} \textquotedblleft{}host\textquotedblright{}, Old High German \emph{wirt}, Proto-Germanic \emph{werduz} \textquotedblleft{}host\textquotedblright{}.
\end{cousinweb}

\begin{grammarlens}[title={What you've built}]
One more word for telling, explaining and describing at B1.
\end{grammarlens}

\subsection*{Your turn}
\begin{itemize}
\raggedright
  \item \emph{Say it:} \emph{bewirten}
  \item \emph{Say it:} \emph{bewirten}, once more
  \item \emph{Say it:} a sentence of your own with \emph{bewirten}
  \item \emph{Recall:} say the German for to unwrap, then the German for the host, then say \emph{bewirten} again
\end{itemize}

\begin{tcolorbox}[breakable,colback=teal!4,colframe=teal!35!black,title={Before you move on}]
What does bewirten mean? (\textquotedblleft{}To treat guests, to serve food and drink to\textquotedblright{}.) Say it once more.
\end{tcolorbox}
\section[der Ehrengast]{der Ehrengast — the guest of honour}

\label{lesson:GE-C307-ehrengast}



\begin{quote}
Before the new one: say the German for the host, then the German for to treat guests.
\end{quote}

\subsection*{You'll want to know: der Ehrengast}
\textbf{der Ehrengast} — \textquotedblleft{}the guest of honour\textquotedblright{}.

\begin{quote}
Der Ehrengast saß neben dem Brautpaar.
\end{quote}

\begin{quote}
\emph{The guest of honour sat next to the bride and groom.}
\end{quote}

Plural: \emph{die Ehrengäste}.

\begin{cousinweb}
\emph{Ehrengast} joins \emph{Ehre} \textquotedblleft{}honour\textquotedblright{}, Old High German \emph{ēra} (EtymDB), and \emph{Gast} \textquotedblleft{}guest\textquotedblright{}, both inherited Germanic words.
\end{cousinweb}

\begin{grammarlens}[title={What you've built}]
One more word for telling, explaining and describing at B1.
\end{grammarlens}

\subsection*{Your turn}
\begin{itemize}
\raggedright
  \item \emph{Say it:} \emph{der Ehrengast}
  \item \emph{Say it:} \emph{der Ehrengast}, once more
  \item \emph{Say it:} a sentence of your own with \emph{der Ehrengast}
  \item \emph{Recall:} say the German for the host, then the German for to treat guests, then say \emph{der Ehrengast} again
\end{itemize}

\begin{tcolorbox}[breakable,colback=teal!4,colframe=teal!35!black,title={Before you move on}]
What does der Ehrengast mean? (\textquotedblleft{}The guest of honour\textquotedblright{}.) Say it once more.
\end{tcolorbox}
\section[die Ansprache]{die Ansprache — the speech, the address}

\label{lesson:GE-C307-ansprache}



\begin{quote}
Before the new one: say the German for to treat guests, then the German for the guest of honour.
\end{quote}

\subsection*{You'll want to know: die Ansprache}
\textbf{die Ansprache} — \textquotedblleft{}the speech, the address\textquotedblright{}.

\begin{quote}
Der Vater der Braut hielt eine kurze Ansprache.
\end{quote}

\begin{quote}
\emph{The bride's father gave a short speech.}
\end{quote}

Plural: \emph{die Ansprachen}. \emph{Eine Ansprache halten}. More personal than \emph{eine Rede}.

\begin{cousinweb}
\emph{Ansprache} is built on \emph{Sprache} \textquotedblleft{}speech\textquotedblright{}, from Old High German \emph{sprāhha}, Proto-Germanic \emph{sprēkijō} (EtymDB), with \emph{an}: words spoken to someone.
\end{cousinweb}

\begin{grammarlens}[title={What you've built}]
One more word for telling, explaining and describing at B1.
\end{grammarlens}

\subsection*{Your turn}
\begin{itemize}
\raggedright
  \item \emph{Say it:} \emph{die Ansprache}
  \item \emph{Say it:} \emph{die Ansprache}, once more
  \item \emph{Say it:} a sentence of your own with \emph{die Ansprache}
  \item \emph{Recall:} say the German for to treat guests, then the German for the guest of honour, then say \emph{die Ansprache} again
\end{itemize}

\begin{tcolorbox}[breakable,colback=teal!4,colframe=teal!35!black,title={Before you move on}]
What does die Ansprache mean? (\textquotedblleft{}The speech, the address\textquotedblright{}.) Say it once more.
\end{tcolorbox}
\section[der Trinkspruch]{der Trinkspruch — the toast (with a drink)}

\label{lesson:GE-C307-trinkspruch}



\begin{quote}
Before the new one: say the German for the guest of honour, then the German for the speech.
\end{quote}

\subsection*{You'll want to know: der Trinkspruch}
\textbf{der Trinkspruch} — \textquotedblleft{}the toast (with a drink)\textquotedblright{}.

\begin{quote}
Nach dem Essen brachte er einen Trinkspruch aus.
\end{quote}

\begin{quote}
\emph{After the meal he proposed a toast.}
\end{quote}

Plural: \emph{die Trinksprüche}. \emph{Einen Trinkspruch ausbringen}.

\begin{cousinweb}
\emph{Trinkspruch} joins \emph{trinken} \textquotedblleft{}to drink\textquotedblright{}, Proto-Germanic \emph{drinkaną}, and \emph{Spruch} \textquotedblleft{}saying\textquotedblright{}, Middle High German \emph{spruch} (EtymDB).
\end{cousinweb}

\begin{grammarlens}[title={What you've built}]
One more word for telling, explaining and describing at B1.
\end{grammarlens}

\subsection*{Your turn}
\begin{itemize}
\raggedright
  \item \emph{Say it:} \emph{der Trinkspruch}
  \item \emph{Say it:} \emph{der Trinkspruch}, once more
  \item \emph{Say it:} a sentence of your own with \emph{der Trinkspruch}
  \item \emph{Recall:} say the German for the guest of honour, then the German for the speech, then say \emph{der Trinkspruch} again
\end{itemize}

\begin{tcolorbox}[breakable,colback=teal!4,colframe=teal!35!black,title={Before you move on}]
What does der Trinkspruch mean? (\textquotedblleft{}The toast (with a drink)\textquotedblright{}.) Say it once more.
\end{tcolorbox}
